\chapter{Background and Technical Foundations}
\label{ch:background}

This chapter reviews the technical concepts that underpin \projname.
Only the ideas that are actually used in the implementation are
covered, at the depth required to make the design decisions in later
chapters legible.

\section{The Machine Learning Workflow}
\label{sec:ml-workflow}

A canonical supervised learning problem is defined as follows. Given a
dataset $\mathcal{D}=\{(x_i, y_i)\}_{i=1}^{n}$ with features
$x_i \in \mathcal{X}$ and targets $y_i \in \mathcal{Y}$ drawn i.i.d.\
from an unknown distribution $\mathbb{P}(X,Y)$, we seek a function
$f:\mathcal{X}\to\mathcal{Y}$ that minimises the expected loss
\begin{equation}
  \mathcal{R}(f) \;=\; \mathbb{E}_{(x,y)\sim\mathbb{P}}\big[\ell(f(x), y)\big],
  \label{eq:risk}
\end{equation}
for a task-appropriate loss $\ell$ \citep{vapnik1998statistical}.
Empirical risk minimisation approximates \eqref{eq:risk} on the
training set,
\begin{equation}
  \hat{f} \;=\; \arg\min_{f\in\mathcal{F}}\;
  \frac{1}{n}\sum_{i=1}^{n} \ell(f(x_i), y_i),
  \label{eq:erm}
\end{equation}
and uses regularisation, cross-validation, and held-out evaluation to
control generalisation error \citep{hastie2009elements}.

An end-to-end applied workflow consists of the stages shown in
Table~\ref{tab:ml-stages}. \projname{} exposes each stage as one or
more tools inside an MCP server.

\begin{table}[H]
\centering
\caption{Canonical stages of a tabular supervised-learning workflow
and the corresponding \projname{} tools.}
\label{tab:ml-stages}
\begin{tabularx}{\linewidth}{@{}lXl@{}}
\toprule
\textbf{Stage} & \textbf{Purpose} & \textbf{Tool in \projname} \\
\midrule
Ingestion & Read file, detect format, report shape/dtypes & \tool{ingest\_dataset} \\
Validation & Enforce schema, detect anomalies & \tool{validate\_schema\_with\_pandera} \\
EDA & Descriptive statistics, correlations, distributions & \tool{run\_full\_eda}, \tool{render\_correlation\_matrix} \\
Splitting \& CV & Train/test split, K-fold CV & inside \tool{run\_parallel\_bake\_off} \\
Training \& selection & Fit candidates, pick champion & \tool{run\_parallel\_bake\_off} \\
Tuning & Hyperparameter search & \tool{trigger\_hyperparameter\_sweep} \\
Explanation & Global feature attributions & \tool{calculate\_shap\_values}, \tool{generate\_feature\_importance\_plot} \\
Export & Reproducible notebook + PDF report & \tool{generate\_jupyter\_notebook}, \tool{compile\_pdf\_report} \\
\bottomrule
\end{tabularx}
\end{table}

\section{Automated Machine Learning}
\label{sec:automl-bg}

AutoML systems automate one or more of the decisions in
Table~\ref{tab:ml-stages}. Formally, an AutoML problem can be phrased
as the joint optimisation of a pipeline $p\in\mathcal{P}$ from a
composite space of preprocessors, estimators and hyperparameters
\citep{feurer2015autosklearn}:
\begin{equation}
  p^{\ast} \;=\; \arg\min_{p\in\mathcal{P}}\;
  \operatorname{CV}\big(\ell,\, p,\, \mathcal{D}_{\text{train}}\big),
  \label{eq:automl-cv}
\end{equation}
where $\operatorname{CV}(\cdot)$ denotes $k$-fold cross-validated risk.
\projname{} implements a small but principled version of
\eqref{eq:automl-cv}: it evaluates four candidate estimators with
$k$-fold cross-validation and selects the one with the best primary
metric.

\section{Data Preprocessing and Feature Engineering}
\label{sec:preproc-bg}

The dominant preprocessing recipe for tabular data
\citep{hastie2009elements, kuhn2019featureengineering} combines
imputation and encoding into a single \code{ColumnTransformer}
\citep{sklearn2011}:
\begin{itemize}[leftmargin=1.3em]
  \item Numerical columns: median imputation, then $z$-score
        standardisation
        $z_j = (x_j - \mu_j)/\sigma_j$.
  \item Categorical columns: most-frequent imputation, then one-hot
        encoding.
\end{itemize}
\projname{} builds exactly this preprocessor in
\file{mcp\_servers/modeling.py} via the \code{\_build\_preprocessor}
function; the same pipeline is embedded into exported Jupyter
notebooks so that a user can re-run it offline.

\section{Exploratory Data Analysis}
\label{sec:eda-bg}

The core of EDA \citep{tukey1977eda} is to summarise, visualise, and
develop intuition for a dataset before modelling. The platform's
\tool{run\_full\_eda} tool produces:

\begin{itemize}[leftmargin=1.3em]
  \item Per-column descriptive statistics for numerical features:
        count, mean, standard deviation, minimum, quartiles, maximum,
        IQR, skewness, and IQR-based outlier count.
  \item Per-column frequency summaries for categorical features:
        count, unique values, top-$k$ frequencies.
  \item Missing-value summaries with per-column counts and
        percentages.
  \item Three plots rendered directly with matplotlib
        \citep{hunter2007matplotlib}: a correlation heatmap, a small
        multiple of histograms, and a bar chart of missing values.
\end{itemize}

Correlation coefficients are computed with the three canonical methods
--- Pearson $\rho$, Spearman $\rho_s$, and Kendall $\tau$ ---
depending on the user's request. For two random variables $X$ and $Y$
with means $\mu_X$, $\mu_Y$ and standard deviations $\sigma_X$,
$\sigma_Y$,
\begin{equation}
  \rho(X, Y) \;=\;
  \frac{\mathbb{E}[(X-\mu_X)(Y-\mu_Y)]}{\sigma_X \sigma_Y}.
  \label{eq:pearson}
\end{equation}

\section{Model Selection and Hyperparameter Optimisation}
\label{sec:hpo-bg}

The candidate models in the \projname{} bake-off are those that have
historically dominated tabular benchmarks
\citep{grinsztajn2022trees, borisov2022deep}: random forests
\citep{breiman2001rf}, gradient-boosted trees (XGBoost
\citep{chen2016xgboost}, LightGBM \citep{ke2017lightgbm}), and simple
linear baselines (logistic regression, ridge regression). Each model
is embedded into an sklearn \code{Pipeline} together with the
preprocessor of Section~\ref{sec:preproc-bg} and evaluated with
$k$-fold cross-validation. The primary metric is F1-weighted for
classification and $R^2$ for regression.

The hyperparameter sweep uses Optuna \citep{akiba2019optuna} with the
Tree-structured Parzen Estimator (TPE) sampler
\citep{bergstra2011algorithms}. TPE models
$p(\lambda \mid \text{score} < \gamma)$ and
$p(\lambda \mid \text{score} \ge \gamma)$ and proposes the point that
maximises the ratio $\ell(\lambda) / g(\lambda)$. \projname{} bounds
the search by wall-clock time, exposing a \code{time\_budget\_seconds}
argument.

\section{Model Evaluation}
\label{sec:eval-bg}

Classification is evaluated with accuracy and the F1 score,
\begin{equation}
  F_1 \;=\; \frac{2\cdot P \cdot R}{P + R},
  \qquad
  P = \frac{\text{TP}}{\text{TP}+\text{FP}},
  \qquad
  R = \frac{\text{TP}}{\text{TP}+\text{FN}},
  \label{eq:f1}
\end{equation}
using the ``weighted'' averaging strategy (per-class F1 weighted by
support) for the multi-class case and ``binary'' otherwise. Regression
uses the root mean squared error (RMSE) and the coefficient of
determination $R^2$,
\begin{equation}
  \text{RMSE} = \sqrt{\frac{1}{n}\sum_{i=1}^{n}(y_i-\hat{y}_i)^2},
  \qquad
  R^2 = 1 - \frac{\sum(y_i-\hat{y}_i)^2}{\sum(y_i-\bar y)^2}.
  \label{eq:rmse-r2}
\end{equation}

\section{Model Explainability}
\label{sec:xai-bg}

The Shapley value assigns to each feature $j$ its expected marginal
contribution to the model output across all coalitions
$S\subseteq\{1,\ldots,d\}\setminus\{j\}$:
\begin{equation}
  \phi_j \;=\;
  \sum_{S\subseteq\{1,\ldots,d\}\setminus\{j\}}
  \frac{|S|!\,(d-|S|-1)!}{d!}
  \big[f(S\cup\{j\}) - f(S)\big].
  \label{eq:shap}
\end{equation}
SHAP \citep{lundberg2017shap} approximates
\eqref{eq:shap} efficiently for tree ensembles (\code{TreeExplainer}
\citep{lundberg2020treeshap}) and generic models
(\code{KernelExplainer}). \projname{} calls \code{TreeExplainer} for
Random Forest, XGBoost, LightGBM, and Gradient Boosting models and
falls back to \code{KernelExplainer} otherwise.

\section{Reproducibility and MLOps Basics}
\label{sec:mlops-bg}

Reproducibility \citep{pineau2021reproducibility} rests on three
pillars: deterministic code, versioned data, and preserved artefacts.
\projname{} addresses the first with fixed \code{random\_state}
arguments plumbed through every tool; the second is the user's
responsibility (raw uploads are preserved verbatim under
\file{data/uploads}); the third is provided by an artefact registry
that assigns each generated file a UUID, a kind, and a metadata
dictionary, and appends it to \file{data/artifacts/\_index.json}.

\section{Large Language Models and Tool Use}
\label{sec:tools-bg}

Modern instruction-tuned LLMs \citep{openai2023gpt4, llama3card,
touvron2023llama} can be prompted with a list of \emph{tool schemas}
--- JSON descriptions of functions the model may call --- and will
emit structured \code{tool\_call} messages when they judge that
external computation is required
\citep{openai2023function, schick2023toolformer}. This turns the LLM
from a text-completion engine into a \emph{planner} that composes
external primitives to satisfy a user goal
\citep{yao2023react, patil2023gorilla}.

\projname{} follows the OpenAI-compatible tool-calling convention as
implemented by the Groq Chat Completions API. The orchestrator
translates the MCP tool catalog into that schema at start-up and
routes each \code{tool\_call} back through the MCP client pool.

\section{The Model Context Protocol}
\label{sec:mcp-bg}

The Model Context Protocol (MCP) \citep{mcpSpec2024, anthropic2024mcp}
is a JSON-RPC-based protocol that standardises how an LLM host
discovers and invokes external capabilities. A conformant server
exposes three kinds of primitives:

\begin{enumerate}[leftmargin=1.4em]
  \item \textbf{Tools} --- typed functions the host may call. Each
        tool has an \code{inputSchema} (JSON Schema) and a natural
        language description.
  \item \textbf{Resources} --- addressable read-only content
        referenced by URI templates.
  \item \textbf{Prompts} --- parameterised prompt templates that the
        host can invoke to shape a user turn.
\end{enumerate}

FastMCP \citep{fastmcp} is a Python implementation of MCP that
generates the required JSON schemas from Python type hints and
docstrings. \projname{} uses FastMCP to define five servers and
communicates with them via the special \code{FastMCPTransport} that
runs the client--server exchange as in-process asynchronous method
calls --- no HTTP, no ports, and no serialisation cost beyond what
the protocol itself requires.

Table~\ref{tab:mcp-primitives} summarises the primitives used by
\projname.

\begin{table}[H]
\centering
\caption{MCP primitives used by \projname.}
\label{tab:mcp-primitives}
\begin{tabularx}{\linewidth}{@{}llXr@{}}
\toprule
\textbf{Primitive} & \textbf{Where} & \textbf{Purpose} & \textbf{Count} \\
\midrule
Tools & All five servers & LLM-callable functions & 11 \\
Prompts & \svc{mcp-eda} & Multi-step template (\tool{eda\_deep\_dive}) & 1 \\
Resources & --- & (Not currently exposed; see Ch.~\ref{ch:conclusion}) & 0 \\
\bottomrule
\end{tabularx}
\end{table}

\section{Streamlit}
\label{sec:streamlit-bg}

Streamlit \citep{streamlit} is a Python library for building
interactive data applications with a re-run-on-input execution model.
Its widgets and state management are well matched to a
research-oriented UI that is developed and maintained by a small team.
\projname{} uses Streamlit only as an interactive shell over the
orchestrator; the UI contains no ML logic and can be swapped for a
CLI, a Jupyter widget layer, or a web front-end without touching
\file{core/}, \file{mcp\_servers/}, or \file{orchestrator/}.
